\chapter[\emj{📋}~DP - Tabulation (Bottom-Up)]{\emj{📋}~DP - Tabulation (Bottom-Up)}

\begin{dsaquees}
Construir una pared de ladrillos 🧱 desde el suelo hacia arriba. 

En lugar de empezar por el techo (el problema grande) y tratar de 
adivinar qué ladrillos necesitas abajo (recursión), empiezas poniendo
los primeros ladrillos en la base (los casos base que ya conoces). 

Una vez que tienes la primera fila, puedes poner la segunda encima 
usando la base como apoyo. Vas llenando tu tabla de resultados paso
a paso hasta que llegas a la altura que querías (el resultado final).
\end{dsaquees}
\begin{dsaotraforma}
Tabulation es la técnica de DP basada en Iteración (Bottom-Up). 

Se llama así porque literalmente llenas una tabla (un array de 1D 
o 2D) con los resultados de los sub-problemas, empezando por los 
más pequeños. Cada nueva entrada de la tabla se calcula usando los
valores que ya están guardados en las posiciones anteriores.

Es generalmente más eficiente que Memoization porque evita el 
"overhead" (sobrecosto) de las llamadas recursivas y no corre el 
riesgo de causar un Stack Overflow.
\end{dsaotraforma}
\begin{dsadiagrama}
Queremos fib(5).
Tabla DP de tamaño 6: [ \_, \_, \_, \_, \_, \_ ]

1. Llenamos bases: dp[0]=0, dp[1]=1
   [ 0, 1, \_, \_, \_, \_ ]

2. Llenamos el resto usando dp[i] = dp[i-1] + dp[i-2]:
   i=2: [ 0, 1, 1, \_, \_, \_ ]
   i=3: [ 0, 1, 1, 2, \_, \_ ]
   i=4: [ 0, 1, 1, 2, 3, \_ ]
   i=5: [ 0, 1, 1, 2, 3, 5 ] 🎉 -> Resultado final
\end{dsadiagrama}
\begin{dsacomofunciona}
1. Inicializa una tabla (array o matriz) con valores default.
2. Define los casos base en las primeras posiciones de la tabla.
3. Escribe un loop (o varios) que recorra la tabla.
4. En cada paso, aplica la lógica para llenar la celda actual 
   usando celdas anteriores.
5. Regresa el valor de la última celda relevante.
\end{dsacomofunciona}
\section*{PSEUDOCÓDIGO}
\hrulefill
\begin{verbatim}
función minMonedasTab(monto, monedas):
    tabla = array de tamaño (monto + 1) con INFINITO
    tabla[0] = 0

    PARA cada moneda EN monedas:
        PARA i desde moneda hasta monto:
            SI tabla[i - moneda] != INFINITO:
                tabla[i] = min(tabla[i], tabla[i - moneda] + 1)

    return tabla[monto]
\end{verbatim}
\section*{CÓDIGO C++}
\hrulefill
\begin{lstlisting}
#include <iostream>
#include <vector>
#include <algorithm>
using namespace std;

// Problema: Climbing Stairs (¿De cuántas formas puedes subir N escalones
// si puedes dar pasos de 1 o 2?) -> O(n) tiempo, O(n) espacio
int climbStairs(int n) {
    if (n <= 2) return n;
    
    vector<int> dp(n + 1);
    dp[1] = 1;
    dp[2] = 2;

    for (int i = 3; i <= n; i++) {
        dp[i] = dp[i - 1] + dp[i - 2];
    }

    return dp[n];
}

// OPTIMIZACIÓN DE ESPACIO: O(n) -> O(1)
// Como solo necesitamos los últimos 2 valores, no hace falta el array.
int climbStairsOptimized(int n) {
    if (n <= 2) return n;
    int prev2 = 1, prev1 = 2;
    for (int i = 3; i <= n; i++) {
        int current = prev1 + prev2;
        prev2 = prev1;
        prev1 = current;
    }
    return prev1;
}

int main() {
    int stairs = 5;
    cout << "Formas de subir " << stairs << " escalones: " 
         << climbStairsOptimized(stairs) << endl; // 8
    return 0;
}
\end{lstlisting}
\begin{dsacomplejidad}
Aspecto        │ Con Tabulation
  ───────────────┼────────────────────────────────────────────
  Tiempo         │ O(Tamaño de la Tabla * Trabajo por Celda)
  Espacio        │ O(Tamaño de la Tabla) -> A veces reducible a O(1)

* Al ser iterativo, el tiempo es muy predecible y constante.
\end{dsacomplejidad}
\begin{dsaentrevistas}

\end{dsaentrevistas}
\begin{dsaresumen}
• Bottom-Up: De lo simple a lo complejo.
• Iterativo: Sin recursión, sin stack overflow.
• Llenas una tabla paso a paso.
• Más eficiente en tiempo y memoria.
• Permite optimizar espacio fácilmente (usando variables en vez de array).
• Es la forma "profesional" de entregar una solución DP.
════════════════════════════════════════════════════════════════
\end{dsaresumen}

